% front_matter.tex -- Abstract and Introduction

\begin{abstract}
As AI agents are deployed in increasingly complex social and economic
settings, measuring their capacity for strategic reasoning, cooperation,
and trust becomes essential for alignment research. We introduce
\textsc{KantBench}, a game-theory benchmark built on the OpenEnv
framework that evaluates AI agents across six canonical games:
Prisoner's Dilemma, Stag Hunt, Hawk-Dove, the Ultimatum Game, the Trust
Game, and a Public Goods Game. Each game is paired with a diverse set of
built-in opponent strategies---from unconditional cooperators to
adaptive and retaliatory policies---enabling controlled measurement of
agent behaviour under varied social pressures. We define six evaluation
metrics that capture cooperation tendency, exploitation resistance,
Pareto efficiency, fairness, adaptability, and an aggregate strategic
reasoning score. By standardizing the game suite, opponent pool, and
evaluation protocol, \textsc{KantBench} provides a reproducible
testbed for studying how language-model agents navigate fundamental
tensions between self-interest and collective welfare.
\end{abstract}

\section{Introduction}
\label{sec:intro}

The deployment of large language model (LLM) agents in real-world
settings---negotiation, resource allocation, market participation---raises
a pressing question: \emph{can we reliably measure and predict how these
agents behave in strategic interactions?} Classical game theory offers a
rich library of settings that distil key social dilemmas into tractable
formal games~\cite{nash1950equilibrium,rapoport1965prisoner,axelrod1984evolution}.
Yet existing AI benchmarks either focus on single-agent decision-making
or embed strategic reasoning inside complex, high-dimensional
environments that make controlled evaluation
difficult~\cite{pan2023machiavelli,leibo2021meltingpot}.

We present \textsc{KantBench}, a lightweight, standardized
benchmark designed to isolate and quantify strategic reasoning in AI
agents. The benchmark consists of six games spanning the core tensions in
social decision-making:

\begin{itemize}[nosep]
    \item \textbf{Matrix games} (Prisoner's Dilemma, Stag Hunt,
          Hawk-Dove) test cooperation, coordination, and conflict
          resolution under simultaneous choice.
    \item \textbf{Sequential games} (Ultimatum, Trust, Public Goods)
          test fairness norms, trust calibration, and free-riding
          incentives.
\end{itemize}

Each game is played against a configurable opponent drawn from a
library of classical strategies (e.g., Tit-for-Tat, Grudger, Pavlov,
Always Defect). This opponent pool lets researchers measure how agents
adapt to cooperative, adversarial, and mixed partners.

\textsc{KantBench} integrates with the OpenEnv framework, exposing
a Gymnasium-compatible API that any reinforcement-learning or
LLM-based agent can consume. A WebSocket client--server architecture
allows remote evaluation, and a built-in tournament runner automates
round-robin matchups with metric aggregation.

Our contributions are:
\begin{enumerate}[nosep]
    \item A curated suite of six game-theory environments with formally
          specified payoff structures (Section~\ref{sec:design}).
    \item A library of seventeen opponent strategies covering
          unconditional, reciprocal, retaliatory, and stochastic play
          (Section~\ref{sec:design}).
    \item An evaluation protocol with six metrics designed to capture
          distinct aspects of strategic and ethical
          reasoning (Section~\ref{sec:evaluation}).
    \item An open-source implementation integrated with the OpenEnv
          framework for reproducible benchmarking
          (Section~\ref{sec:architecture}).
\end{enumerate}
